\documentclass[10pt,a4paper]{amsart}
\usepackage[margin=19mm]{geometry}
\usepackage{amsmath,amssymb,mathtools}
\usepackage[scr=rsfs,cal=euler]{mathalfa}
\usepackage{xfrac}
\usepackage[unicode,hidelinks,bookmarks=false]{hyperref}
\newcommand{\ie}{i.e.\ }
\newcommand{\cf}{cf.\ }
\newcommand{\resp}{respectively\ }
\newcommand{\Subsection}{\subsection}
\providecommand{\nobreakdash}{\nobreak\hbox{-}}
\setlength{\emergencystretch}{2em}
\multlinegap0pt
\usepackage{amssymb}
\usepackage[shortlabels]{enumitem}
\usepackage{caption}
\usepackage{booktabs}
\usepackage{float}
\usepackage{xcolor}
\usepackage[normalem]{ulem}
\newtheorem{thm}{Theorem}
\newtheorem{lemma}{Lemma}
\newcommand\dd{{\mathrm d}}
\newcommand{\f}[2]{\frac{#1}{#2}}
\title{Expander estimates for cubes}
\author{J\"org Brüdern and Simon L. Rydin Myerson}
\date{Source: arXiv:2409.16795v2 (CC BY 4.0)}
\begin{document}
\maketitle

\noindent\emph{Source and licence.} Expander estimates for cubes. Version arXiv:2409.16795v2 (CC BY 4.0), distributed by arXiv under the Creative Commons Attribution 4.0 International licence, http://creativecommons.org/licenses/by/4.0/, as recorded in the arXiv licence metadata for this exact version and shown on the abstract page https://arxiv.org/abs/2409.16795v2. Full licence notice: \texttt{licenses/arxiv-2409.16795-CC-BY-4.0.txt}.\par\smallskip

\noindent\textit{Editorial scope.} Counted unit: the article's lemma L34 (source0.tex lines 469--472). The article's complete original proof of it is delivered with it (source0.tex lines 474--498, one \texttt{proof} environment of 25 lines). No mathematics is written, repaired or re-derived: the statement and the proof are the article's own lines, in the article's own notation, and no retained line is substituted or re-flowed.  Retained uncounted as the source-local context the proof needs: lines 358--360; lines 375--375; lines 457--457; lines 462--462; lines 465--466.  Nothing was omitted from any retained span, so the package carries no omission record.  No retained line contains a citation, so the wrapper carries no bibliography and no bibliography entry.  No retained line contains a figure environment, an image-inclusion command or drawing code, so no image is delivered and the package declares no asset dependency.  Editorial formatting only, in addition: an AMS \texttt{amsart} preamble that restates the article's own declarations and macros used by the retained lines, namely the environments \texttt{thm}, \texttt{lemma} on separate counters, together with the standard packages those lines need.  The retained environments print in this document as Theorem 1, Lemma 1 in order of appearance, whereas the article prints the same items with the numbering of its own full document.  The complete native source of the article as distributed in the arXiv e-print for this exact version is bundled unchanged as \texttt{sources/165788/source0.tex}.
\begin{thm}\label{quadWeyl} Let $a_1,a_2\in\mathbb Z$, $q\in\mathbb N$, $\alpha_j\in\mathbb R$ and $\beta_j = \alpha_j-a_j/q$. Suppose that $|\beta_1|\le 1/(2q)$. Then
$$ f(\alpha_1,\alpha_2) - q^{-1}S(q,a_1,a_2)I(\beta_1,\beta_2) \ll (q,a_2)^{1/2} (q+qX^2|\beta_2|)^{1/2} \log q. $$
\end{thm}

\begin{lemma}\label{L23} If $(a_1,a_2,q)=1$, then $S(q,a_1,a_2)\ll q^{1/2}$. If in addition $(q,a_2)>1$, then $S(q,a_1,a_2)=0$. \end{lemma}

\begin{equation}\label{maj} 0\le a\le q\le P \quad \text{ and } \quad\alpha = a/q +\beta, \quad q|\beta|\le (6HP)^{-1}. \end{equation}

\begin{equation}\label{37} g(\alpha_1,\alpha_2;X) = \sum_{X<x\le 2X} e(\alpha_1 x+\alpha_2 x^2), \quad J(\beta_1,\beta_2;X) = \int_X^{2X} e(\beta_1 u+\beta_2 u^2)\,\dd u, \end{equation}

\begin{equation}\label{38} F_w(\alpha) = \sum_{d\mid w} \mu(d) \sum_{h\le H/d} e(\alpha d^3h^3)
g(3\alpha d^3h^2, 3\alpha d^3h; P/d). \end{equation}

\begin{lemma}\label{L34} Let $\alpha\in\mathfrak M$  and suppose that \eqref{maj} applies. Then, uniformly in $w$, one has
$$ F_w(\alpha) = \sum_{d\mid w}\mu(d) \sum_{h\le H/d} e(\alpha d^3 h^3)
\frac{S(q,3ad^3h^2,3ad^3h)}{qd} J(3h^2d^2\beta, 3hd\beta) + O(P^{1+\varepsilon}).$$ 
\end{lemma}

\begin{proof}
For $j=1$ and $2$ we have
$$ 3\alpha d^3 h^j = \frac{3ad^3h^j}{q} + 3h^jd^3\beta. $$
Here we cancel  common factors between $q$ and $d^3$. On writing $q_0=q/(q,d^3)$ and $d_0=d^3/(q,d^3)$, we have $(q_0,d_0)=1$ and
$$    3\alpha d^3 h^j = \frac{3ah^jd_0}{q_0} + 3h^jd^3\beta \quad (j=1,2). $$
We wish to apply Theorem \ref{quadWeyl} with $X=2P/d$ and $X=P/d$ to evaluate the sum $g(3\alpha d^3h^2, 3\alpha d^3h; P/d)$. We should then choose 
$$\alpha_1= 
3\alpha d^3 h^2, \quad \alpha_2 =3\alpha d^3 h, \quad \beta_1=3 d^3 h^2 \beta\quad \beta_2 = 3 d^3 h\beta. $$
Recalling that $hd\le H$, we see that $|\beta_1|\le 3H^3/(6HPq)=1/(2q)$. Theorem \ref{quadWeyl} is therefore applicable in the way we desired, and we infer that
\begin{equation}\label{39} 
  g(3\alpha d^3h^2, 3\alpha d^3h; P/d) = q_0^{-1} S(q_0, 3ah^2d_0, 3ahd_0) J(3 d^3 h^2 \beta,3 d^3 h\beta;P/d) + E \end{equation}
where
$$ E \ll (q_0,h)^{1/2} (q_0 +q_0 P^2dh |\beta|)^{1/2} \log q. $$
But $q_0\mid q$,  $q \le P$ and $|\beta|\le 1/(6HPq)$ simplify this bound to
$ E\ll (q,h)^{1/2} P^{1/2+\varepsilon}$. If one inserts the expansion \eqref{39} to \eqref{38}, the errors $E$ sum to an amount no larger than
$$ P^{1/2+\varepsilon} \sum_{d\le H} \sum _{h\le H/d} (q,h)^{1/2} \ll P^{1+2\varepsilon}. $$
We have now proved that
\begin{multline*}F_w(\alpha) = \\
\sum_{d\mid w}\mu(d) \sum_{h\le H/d} e(\alpha d^3 h^3) \f{S(q_0, 3ah^2d_0, 3ahd_0)}{q_0}J(3 d^3 h^2 \beta,3 d^3 h\beta;P/d) +O(P^{1+\varepsilon}).\end{multline*}
By Lemma \ref{L23}, we have 
$$ q_0^{-1} S(q_0, 3ah^2d_0, 3ahd_0)= q^{-1} S(q,3ad^3h^2, 3ad^3h). $$
Further, the substitution $u=v/d$ in \eqref{37} gives 
$$ J(3 d^3 h^2 \beta,3 d^3 h\beta;P/d) = d^{-1} J(3h^2d^2\beta, 3hd\beta). $$
Collecting all this together, one arrives at the claim of the lemma.
\end{proof} 

\end{document}
